\section{Conclusions}\label{sec:conclusions}

We built a learned warm start for phase-fixed weighted Gerchberg--Saxton: a permutation-equivariant set network with \nParams{} parameters that predicts one initial phase per trap, trained without labels through an unrolled copy of the iteration, with an acceptance test fixed in advance. On $\nGeo$ held-out geometries it does not meet the test. The network start with an early phase freeze reaches a relative standard deviation of 1 percent in a median of \itNThree{} iterations (95\% interval \itNThreeLo{} to \itNThreeHi) against \itRTwelve{} for the random start of the published method, a factor of \ratioRoverN{} (interval \ratioRoverNLo{} to \ratioRoverNHi) where two was required,
and its final efficiency is \effdNThree{} percent from the published method's where 1 percent was allowed.

The four-arm ablation attributes the gain to the freeze index and not to the network: a random start with the same early freeze needs \itRThree{} iterations, the network at the published freeze index needs \itNTwelve{}, and a start from the phases of a converged run saves only \OTwelveGain{} iteration at that index. The early freeze has a price that the median hides: \neverRThree{} percent of the runs never reach 1 percent within \budget{} iterations, against \neverRTwelve{} percent for the published default, and the final efficiency is about \effdRThree{} percent lower. The accounting bound on the weights holds on every run and is loose, so the restart of the weights does not explain the iteration count; the freeze index does.

The model is not released. The questions that remain are different designs, not tunings: a start trained on labels where converged phases are determined, or a change of the freeze rule that lets a start carry information through the iteration. The results hold on a scalar Fourier model of the device and for one trained network; the experiment that would extend them is a measurement on an instrument with its device model and camera feedback, and a study of the variation across training seeds.
